\section{Common-task errors and causal resources}\label{sec:composition}
\begin{proposition}[Uniform perturbation of an executed observer]\label{prop:perturbation}
Let $P$ be a common-grid observer map on $C_K^m$ satisfying \eqref{eq:ensemble-minor}. Let $\widetilde P$ be another positive mass-preserving ensemble map, with
\begin{equation}\label{eq:operator-error}
 \sup_{Z\in C_K^m,\ e_{\rm tot}(Z)=1}
       \norm{(\widetilde P-P)Z}_{\oplus K}\le\tau.
\end{equation}
Writing $z^*$ for the invariant ensemble of $P$, every invariant normalized ensemble $\widetilde z$ of $\widetilde P$ obeys
\begin{equation}\label{eq:perturbation}
 \norm{\widetilde z-z^*}_{\oplus K}\le b\tau/\gamma.
\end{equation}
From the same initial ensemble in $C_K^m$, the distance between their time-$t$ ensembles is at most $b\tau/\gamma$ for every $t\ge0$. The latter statement also holds for time-dependent perturbed maps having the same per-call bound. If $b\tau<\gamma$, the homogeneous $\widetilde P$ itself has a unique invariant ensemble and block contraction coefficient at most $1-\gamma+b\tau$.

For fixed bounded audit readouts, the stationary and finite-average risk perturbations are at most $L_0b\tau/\gamma+e_{\rm out}$, where $e_{\rm out}$ is a separately certified expected loss change in the readout. The coefficient $b/\gamma$ does not depend on $M$.
\end{proposition}
\begin{proof}
For zero-mass differences, $P^j$ contracts by at most
$(1-\gamma)^{\lfloor j/b\rfloor}$. Along a perturbed execution the one-call residual has mass zero and norm at most $\tau$. Expanding the difference by inserting one residual at a time therefore gives
\[
 \norm{\widetilde Z_t-Z_t}_{\oplus K}
 \le\tau\sum_{j=0}^{t-1}(1-\gamma)^{\lfloor j/b\rfloor}
 \le b\tau/\gamma.
\]
This uses no stationary calibration of the perturbed execution. An invariant ensemble exists by compactness and continuity of its finite-dimensional positive map. Start the comparison there and let $t\to\infty$ to prove \eqref{eq:perturbation}.

For any difference $D$ of total mass zero, a positive/negative base-norm decomposition and \eqref{eq:operator-error} give
$\norm{(\widetilde P-P)D}\le\tau\norm D$. Both maps are contractions on their full base-norm spaces. Telescoping the $b$th powers bounds their operator difference on zero-mass vectors by $b\tau$. Hence $\widetilde P^b$ contracts that space by at most $1-\gamma+b\tau$. Its fixed point is unique and is fixed by $\widetilde P$, as before. The loss is a bounded affine functional on normalized ensembles, proving the risk statement.
\end{proof}

\begin{proposition}[Typed causal simulation]\label{prop:morphism}
Fix the same protected stopping request, response space and loss in $[0,L]$. A parameter-independent causal simulator from a source experiment to the target must provide a legal strategy lift and action map, causal report transformation, $R$ counted internal labels, a call and physical-time map, and compatible task/audit readouts. If its complete stopped and scored law differs by at most $\varepsilon$ in total variation, uniformly over the target policies compared, an $M$-label target observer is implemented with at most $RM$ source labels and risk at most its target risk plus $L\varepsilon+e_{\rm out}$.

Compatible serial simulators multiply label budgets, compose call maps and add complete-law errors. A reverse lower requires a reverse simulator with its own resource map. A per-call full-law error alone gives a finite-prefix bound proportional to the number of calls; the uniform ensemble conclusion of Proposition~\ref{prop:perturbation} requires its stronger common-ensemble hypotheses.
\end{proposition}
\begin{proof}
Store the simulator and observer labels as a pair. Route actions through the legal lift and generate each target report without using its future or an unknown true parameter. The observer's fresh private coins are included in the complete joint law. Compatibility with the external request fixes the same scored task. For total variation defined by a supremum of event differences, integrating a loss in $[0,L]$ changes expectation by at most $L\varepsilon$. Add the separately certified readout perturbation. Causal composition and the triangle inequality prove the serial claim. This standard coupling argument is not an assertion of infinite-transcript total-variation convergence, LDP equivalence or optimality of the product state count.
\end{proof}

\begin{corollary}[An attained joint resource--risk law]\label{cor:joint}
There is a physical extension of every experiment in Theorem~\ref{thm:main}, with one summed-square protected task, an unobservable calibration alternative $\pm\delta$ and an attained bit-report deficiency $\varepsilon\in[0,1/2]$, such that for all sufficiently large $M$,
\begin{equation}\label{eq:joint}
 \begin{split}
 cM^{-2/d}+\delta^2+2\varepsilon
 &\le\calR^{\rm ext}_{\rm hom}(n,M)-J_*(n)\\
 &\le C\bigl(M^{-2/d}+n^{-1}\bigr)+\delta^2+2\varepsilon.
 \end{split}
\end{equation}
The pathwise maximum physical-call budget is $N=n+3$. On the matching region $n\ge c_*M^{2/d}$, this gives
$\calR^{\rm ext}_{\rm hom}-J_*(N-3)\asymp M^{-2/d}+\delta^2+\varepsilon$.
The weighted profile $\Xi(M)$ may replace $M^{-2/d}$ under Theorem~\ref{thm:strata}. Numerical or model error ceilings are not automatically lower bounds.
\end{corollary}
\begin{proof}
When the hidden audit request arrives, stop all core controls and hold the core state unchanged in the explicitly enlarged post-request interface. Only then prepare an independent fair bit $B\in\{-1,1\}$ and make one paid report call which reveals it with probability $1-a$ and otherwise returns a marked erasure. Commit the response before a final audit reveals the core audit and $B$. Bit responses lie in $[-1,1]$.

On the two prescribed bit-preparation parameters, the deficiency to an exact bit report completed before the audit is exactly $a/2$. Filling an erasure with a fair bit attains the upper. For any simulator the sum of the two error probabilities on an erased report is at least $a$, proving the lower. The fair-bit Bayes squared risk is exactly $a$, obtained with readouts $-1,0,1$. Set $a=2\varepsilon$.

An independent calibration coordinate is fixed at $+\delta$ or $-\delta$ and invisible in all pre-audit reports. Averaging these two alternatives gives minimax squared risk at least $\delta^2$, attained by output zero. Add this coordinate, the bit loss and the core loss. Since the bit is prepared after controls stop, it cannot have been stored by earlier control actions in the physical system. Its report conveys no core information conditional on the previous observer label. Arbitrary randomization and the extra terminal labels therefore do not evade the core finite-$n$ lower. Independence and the product prior give the additive lower in \eqref{eq:joint}.

For the upper use $m$ running core labels and three post-request readout copies for each, storing the bit posterior $-1,0,1$. At most $4m+2$ labels implement all roles. Taking $m$ a fixed fraction of $M$ preserves the geometric order and, because the core mixing constant is independent of $m$, also preserves the $C/n$ entrance bound. This is a sufficient implementation, not a minimal-state claim for the enlarged alphabet. Preparation, at most $n$ core reports, the bit report and the single joint audit total at most $n+3$ calls. The physical post-request hold and audit timing are part of this constructed extension, not assumed freely available in an unrelated experiment.
\end{proof}
